% Options for packages loaded elsewhere
\PassOptionsToPackage{unicode}{hyperref}
\PassOptionsToPackage{hyphens}{url}
\documentclass[
]{article}
\usepackage{xcolor}
\usepackage{amsmath,amssymb}
\setcounter{secnumdepth}{-\maxdimen} % remove section numbering
\usepackage{iftex}
\ifPDFTeX
  \usepackage[T1]{fontenc}
  \usepackage[utf8]{inputenc}
  \usepackage{textcomp} % provide euro and other symbols
\else % if luatex or xetex
  \usepackage{unicode-math} % this also loads fontspec
  \defaultfontfeatures{Scale=MatchLowercase}
  \defaultfontfeatures[\rmfamily]{Ligatures=TeX,Scale=1}
\fi
\usepackage{lmodern}
\ifPDFTeX\else
  % xetex/luatex font selection
\fi
% Use upquote if available, for straight quotes in verbatim environments
\IfFileExists{upquote.sty}{\usepackage{upquote}}{}
\IfFileExists{microtype.sty}{% use microtype if available
  \usepackage[]{microtype}
  \UseMicrotypeSet[protrusion]{basicmath} % disable protrusion for tt fonts
}{}
\makeatletter
\@ifundefined{KOMAClassName}{% if non-KOMA class
  \IfFileExists{parskip.sty}{%
    \usepackage{parskip}
  }{% else
    \setlength{\parindent}{0pt}
    \setlength{\parskip}{6pt plus 2pt minus 1pt}}
}{% if KOMA class
  \KOMAoptions{parskip=half}}
\makeatother
\usepackage{color}
\usepackage{fancyvrb}
\newcommand{\VerbBar}{|}
\newcommand{\VERB}{\Verb[commandchars=\\\{\}]}
\DefineVerbatimEnvironment{Highlighting}{Verbatim}{commandchars=\\\{\}}
% Add ',fontsize=\small' for more characters per line
\newenvironment{Shaded}{}{}
\newcommand{\AlertTok}[1]{\textcolor[rgb]{1.00,0.00,0.00}{\textbf{#1}}}
\newcommand{\AnnotationTok}[1]{\textcolor[rgb]{0.38,0.63,0.69}{\textbf{\textit{#1}}}}
\newcommand{\AttributeTok}[1]{\textcolor[rgb]{0.49,0.56,0.16}{#1}}
\newcommand{\BaseNTok}[1]{\textcolor[rgb]{0.25,0.63,0.44}{#1}}
\newcommand{\BuiltInTok}[1]{\textcolor[rgb]{0.00,0.50,0.00}{#1}}
\newcommand{\CharTok}[1]{\textcolor[rgb]{0.25,0.44,0.63}{#1}}
\newcommand{\CommentTok}[1]{\textcolor[rgb]{0.38,0.63,0.69}{\textit{#1}}}
\newcommand{\CommentVarTok}[1]{\textcolor[rgb]{0.38,0.63,0.69}{\textbf{\textit{#1}}}}
\newcommand{\ConstantTok}[1]{\textcolor[rgb]{0.53,0.00,0.00}{#1}}
\newcommand{\ControlFlowTok}[1]{\textcolor[rgb]{0.00,0.44,0.13}{\textbf{#1}}}
\newcommand{\DataTypeTok}[1]{\textcolor[rgb]{0.56,0.13,0.00}{#1}}
\newcommand{\DecValTok}[1]{\textcolor[rgb]{0.25,0.63,0.44}{#1}}
\newcommand{\DocumentationTok}[1]{\textcolor[rgb]{0.73,0.13,0.13}{\textit{#1}}}
\newcommand{\ErrorTok}[1]{\textcolor[rgb]{1.00,0.00,0.00}{\textbf{#1}}}
\newcommand{\ExtensionTok}[1]{#1}
\newcommand{\FloatTok}[1]{\textcolor[rgb]{0.25,0.63,0.44}{#1}}
\newcommand{\FunctionTok}[1]{\textcolor[rgb]{0.02,0.16,0.49}{#1}}
\newcommand{\ImportTok}[1]{\textcolor[rgb]{0.00,0.50,0.00}{\textbf{#1}}}
\newcommand{\InformationTok}[1]{\textcolor[rgb]{0.38,0.63,0.69}{\textbf{\textit{#1}}}}
\newcommand{\KeywordTok}[1]{\textcolor[rgb]{0.00,0.44,0.13}{\textbf{#1}}}
\newcommand{\NormalTok}[1]{#1}
\newcommand{\OperatorTok}[1]{\textcolor[rgb]{0.40,0.40,0.40}{#1}}
\newcommand{\OtherTok}[1]{\textcolor[rgb]{0.00,0.44,0.13}{#1}}
\newcommand{\PreprocessorTok}[1]{\textcolor[rgb]{0.74,0.48,0.00}{#1}}
\newcommand{\RegionMarkerTok}[1]{#1}
\newcommand{\SpecialCharTok}[1]{\textcolor[rgb]{0.25,0.44,0.63}{#1}}
\newcommand{\SpecialStringTok}[1]{\textcolor[rgb]{0.73,0.40,0.53}{#1}}
\newcommand{\StringTok}[1]{\textcolor[rgb]{0.25,0.44,0.63}{#1}}
\newcommand{\VariableTok}[1]{\textcolor[rgb]{0.10,0.09,0.49}{#1}}
\newcommand{\VerbatimStringTok}[1]{\textcolor[rgb]{0.25,0.44,0.63}{#1}}
\newcommand{\WarningTok}[1]{\textcolor[rgb]{0.38,0.63,0.69}{\textbf{\textit{#1}}}}
\setlength{\emergencystretch}{3em} % prevent overfull lines
\providecommand{\tightlist}{%
  \setlength{\itemsep}{0pt}\setlength{\parskip}{0pt}}
\usepackage{bookmark}
\IfFileExists{xurl.sty}{\usepackage{xurl}}{} % add URL line breaks if available
\urlstyle{same}
\hypersetup{
  hidelinks,
  pdfcreator={LaTeX via pandoc}}

\author{}
\date{}

\begin{document}

\section{AEP-Media: Reusable Research Software for Offline Validation of
Time-Aware Media Evidence
Bundles}\label{aep-media-reusable-research-software-for-offline-validation-of-time-aware-media-evidence-bundles}

\subsection{Authors}\label{authors}

Bin Zhang

\subsection{Abstract}\label{abstract}

AEP-Media is reusable research software for representing a media-bearing
operation as a locally checkable evidence bundle. It provides a minimal
time-aware media evidence profile, JSON schemas, controlled examples,
validators, command-line tools, offline bundle build and verification,
strict declared time-trace validation, adapter-only ingestion
interfaces, evaluation matrices, and release packaging reports. The
software is intended for researchers studying operation accountability,
media evidence packaging, provenance boundaries, and reproducible
validation artifacts. AEP-Media checks profile identity, required
fields, reference closure, artifact hashes, path safety, bundle
checksums, declared clock-trace coverage, clock-trace summaries, and
offset/jitter thresholds. Its evaluation materials exercise valid,
invalid, tamper, adapter-ingestion, and optional-tool reporting cases
with expected pass/fail behavior. The current claim is local validation
and fixture-based adapter ingestion. AEP-Media does not claim legal
admissibility, non-repudiation, trusted timestamping, real PTP proof,
full MP4 PRFT parsing, real C2PA signature verification, chain of
custody, or production deployment.

\subsection{Keywords}\label{keywords}

media evidence; operation accountability; research software; provenance;
offline validation; evidence bundle; time trace

\subsection{Software Availability}\label{software-availability}

\begin{itemize}
\tightlist
\item
  Software name: AEP-Media, within \texttt{agent-evidence}
\item
  Repository: \texttt{https://github.com/joy7758/agent-evidence}
\item
  License: Apache-2.0
\item
  Primary package: \texttt{agent\_evidence}
\item
  Programming language: Python
\item
  Documentation: \texttt{spec/}, \texttt{schema/},
  \texttt{examples/media/}, \texttt{demo/}, and \texttt{docs/reports/}
\item
  Tests: media profile, bundle, strict-time, adapter, evaluation, and
  release-pack tests under \texttt{tests/}
\item
  Archive: AEP-Media-specific archive DOI pending release archive;
  action required before final submission
\end{itemize}

\subsection{Statement of Need}\label{statement-of-need}

Researchers working with media-bearing operations often need to review
several disconnected artifacts: media files, logs, provenance
declarations, timing metadata, policy references, and validation
reports. These objects are difficult to inspect consistently when they
are not bound into a local evidence object with explicit validation
semantics. AEP-Media addresses this gap by providing a reusable software
layer for turning such materials into locally checkable evidence bundles
with machine-readable failure codes.

The problem is deliberately bounded. AEP-Media is not a legal evidence
platform and does not attempt to prove external authenticity. Its
purpose is to make the local validation boundary explicit: which fields
must be present, which references must close, which files must match
declared hashes, which time traces cover a declared window, and which
failures are reported when a controlled rule is broken. This local
conformance layer is useful before stronger external assurance
mechanisms such as signing, trusted timestamping, or deployed capture
pipelines are added.

\subsection{Software Description}\label{software-description}

AEP-Media extends the \texttt{agent-evidence} package with a
media-focused profile and validation path. A media evidence statement
binds an actor, subject, operation, policy, constraints, time context,
media artifacts, provenance references, evidence notes, and validation
expectations into one object. The profile validator checks that the
statement is structurally complete and locally consistent.

The offline bundle layer packages a statement and its artifacts into a
portable directory. The builder copies artifacts, rewrites paths to
bundle-local locations, recomputes hashes and sizes, writes checksums,
and emits validation reports. The verifier rejects unsafe paths, missing
artifacts, checksum mismatches, and media profile failures. This allows
review without relying on the original runtime layout.

The strict-time layer validates declared clock-trace artifacts. It
checks trace references, artifact roles, trace hashes, trace profile
identity, collection window coverage, sample validity, summary
recomputation, offset thresholds, jitter thresholds, and media binding
to the declared time context.

The adapter layer is ingestion-only. LinuxPTP-style logs, FFmpeg
PRFT-style timing metadata, and C2PA-like manifest metadata can be
normalized into AEP-Media reports. Adapter reports distinguish fixture
ingestion from optional external-tool reporting and do not imply
external verification unless a future environment-specific run records
it.

\subsection{Architecture and
Functionality}\label{architecture-and-functionality}

The main modules are:

\begin{itemize}
\tightlist
\item
  \texttt{agent\_evidence/media\_profile.py}: media profile validation
\item
  \texttt{agent\_evidence/media\_bundle.py}: offline bundle build and
  verification
\item
  \texttt{agent\_evidence/media\_time.py}: strict declared time-trace
  validation
\item
  \texttt{agent\_evidence/adapters/linuxptp.py}: LinuxPTP-style trace
  ingestion
\item
  \texttt{agent\_evidence/adapters/ffmpeg\_prft.py}: FFmpeg PRFT-style
  metadata ingestion
\item
  \texttt{agent\_evidence/adapters/c2pa\_manifest.py}: C2PA-like
  manifest ingestion
\item
  \texttt{agent\_evidence/media\_evaluation.py}: default evaluation
  matrix
\item
  \texttt{agent\_evidence/media\_adapter\_evaluation.py}:
  adapter-inclusive evaluation
\item
  \texttt{agent\_evidence/media\_optional\_tools.py}: optional
  external-tool reporting
\item
  \texttt{agent\_evidence/media\_release\_pack.py}: release pack
  assembly
\end{itemize}

Typical commands are:

\begin{Shaded}
\begin{Highlighting}[]
\ExtensionTok{agent{-}evidence}\NormalTok{ validate{-}media{-}profile examples/media/minimal{-}valid{-}media{-}evidence.json}
\ExtensionTok{agent{-}evidence}\NormalTok{ build{-}media{-}bundle examples/media/minimal{-}valid{-}media{-}evidence.json }\AttributeTok{{-}{-}out}\NormalTok{ /tmp/aep{-}media{-}bundle{-}check}
\ExtensionTok{agent{-}evidence}\NormalTok{ verify{-}media{-}bundle /tmp/aep{-}media{-}bundle{-}check}
\ExtensionTok{agent{-}evidence}\NormalTok{ validate{-}media{-}time{-}profile examples/media/time/minimal{-}valid{-}time{-}aware{-}media{-}evidence.json}
\ExtensionTok{agent{-}evidence}\NormalTok{ run{-}media{-}evaluation }\AttributeTok{{-}{-}out}\NormalTok{ demo/output/media\_evaluation\_demo}
\end{Highlighting}
\end{Shaded}

The repository includes valid examples, controlled invalid examples,
strict-time fixtures, adapter fixtures, and tamper cases. These allow
researchers to reproduce the intended pass/fail behavior without
installing LinuxPTP, FFmpeg, ffprobe, or C2PA.

\subsection{Reproducibility and
Evaluation}\label{reproducibility-and-evaluation}

AEP-Media is evaluated as a conformance-oriented software artifact. The
evaluation targets reproducibility, diagnosability, and validation
coverage rather than throughput or field deployment. Prior release
reports record:

\begin{itemize}
\tightlist
\item
  default evaluation: 18 cases, \texttt{unexpected=0};
\item
  adapter-inclusive evaluation: 26 cases, \texttt{unexpected=0};
\item
  optional-tool reporting evaluation: 23 cases, \texttt{unexpected=0};
\item
  combined adapter and optional-tool evaluation: 31 cases,
  \texttt{unexpected=0}.
\end{itemize}

The case categories include valid conformance cases, invalid single-rule
cases, bundle tamper cases, strict-time failures, adapter-ingestion
cases, and optional-tool reporting cases. Expected failures include
missing time context, media hash mismatch, unresolved policy reference,
bundle checksum mismatch, path escape, missing clock-trace reference,
clock offset threshold exceedance, clock window mismatch, missing PRFT
metadata, and declared invalid C2PA-like signature status.

\subsection{Impact}\label{impact}

AEP-Media supports research on operation accountability, media evidence
packaging, provenance boundaries, and reproducible validation artifacts.
It gives researchers a concrete software artifact for asking whether a
declared media evidence package is internally consistent and portable.
The software may also serve as a baseline for future work that adds real
external tool runs, signed manifests, trusted timestamps, or deployed
capture pipelines.

The impact is not that AEP-Media proves authenticity. Its value is that
it makes the local evidence object explicit and reproducibly checkable,
which is a prerequisite for stronger assurance layers.

\subsection{Limitations and Claim
Boundary}\label{limitations-and-claim-boundary}

AEP-Media detects inconsistencies inside a declared evidence package. It
does not establish that the original media capture event was truthful,
authorized, or unmodified before packaging. It does not provide legal
admissibility, non-repudiation, trusted timestamping, real PTP proof,
full MP4 PRFT parsing, real C2PA signature verification, chain of
custody, production deployment, or broad forensic sufficiency.

Optional external-tool paths are reporting paths. The reproducible
baseline remains local validation and fixture-based adapter ingestion.

\subsection{Future Work}\label{future-work}

Future work should keep the local profile semantics stable while adding
environment-specific evidence appendices: LinuxPTP traces captured on
equipped Linux hosts, ffprobe output from PRFT-bearing media, C2PA CLI
reports for real signed assets, independent reproduction of the
evaluation matrix, and an AEP-Media-specific archived release DOI.

\subsection{Declaration of Generative AI and AI-assisted Technologies in
the Manuscript Preparation
Process}\label{declaration-of-generative-ai-and-ai-assisted-technologies-in-the-manuscript-preparation-process}

During the preparation of this work, the author used OpenAI
ChatGPT/Codex for manuscript organization, implementation scaffolding,
command generation, and wording refinement. After using these tools, the
author reviewed and edited the content as needed and takes full
responsibility for the content of the submitted article.

\subsection{References}\label{references}

{[}1{]} W3C, ``PROV Overview,'' World Wide Web Consortium.

{[}2{]} W3C, ``PROV-DM: The PROV Data Model,'' World Wide Web
Consortium.

{[}3{]} Coalition for Content Provenance and Authenticity, ``C2PA
Technical Specification.''

{[}4{]} FFmpeg Project, ``ffprobe Documentation.''

{[}5{]} LinuxPTP Project, ``ptp4l and phc2sys Documentation.''

{[}6{]} GStreamer Project, ``GstPtpClock Documentation.''

{[}7{]} JSON Schema, ``JSON Schema Core and Validation Specifications.''

{[}8{]} Association for Computing Machinery, ``Artifact Review and
Badging.''

{[}9{]} in-toto Project, ``in-toto: Providing farm-to-table guarantees
for bits and bytes.''

{[}10{]} Supply-chain Levels for Software Artifacts, ``SLSA
Provenance.''

{[}11{]} R. Kahn and R. Wilensky, ``A Framework for Distributed Digital
Object Services.''

{[}12{]} DONA Foundation, ``Digital Object Interface Protocol
Specification.''

\end{document}
